\subsection{Sampling Methods}

\content{
        \begin{definition} A sampling method is called \textbf{biased} is every member
        of the population doesn't have equal likelihood of being in the sample.
        \end{definition}
}{% presentation
}{% nopresentation
}{% comments
}

\content{
        \begin{definition} A \textbf{simple random sample} is when you list 
        the entire population, and randomly select (hopefully by computer) a
        fixed number of them.
        \end{definition}
}{% presentation
}{% nopresentation
}{% comments
}

\content{
        \begin{definition} The natural variation of samples is called 
        \textbf{sampling variability}. This is unavoidable, and usually is not
        an issue.
        \end{definition}
}{% presentation
}{% nopresentation
}{% comments
}

\content{
        \begin{definition} In \textbf{stratified sampling}, a
        population is divided into a number of subgroups (or strata).
        Random samples are then taken from each strata with samples
        sizes proportional to the size of the strata in the
        population.
        \end{definition}
}{% presentation
\vspace{2in}
}{% nopresentation
\vspace{1in}
}{% comments
}

\content{
        \begin{example}
        Suppose in a particular state that previous data indicated
        that the electorate was comprised of  39\% Democrats, 37\%
        Republicans and 24\% independents. If a sample of 1000 people
        is desired, describe how stratified sampling can be used to
        obtain a good sample.
        \end{example}
}{% presentation
\vspace{2in}
}{% nopresentation
\vspace{1in}
}{% comments
}

\content{
        \begin{definition}
        Quota sampling is a variation on stratified sampling, wherein
        samples are collected in  each subgroup until the desired
        quota is met.
        \end{definition}
        \begin{example}
        Suppose the pollsters call people at random, but once they
        have met their quota of 390  Democrats, they only gather
        people who do not identify themselves as a Democrat.
        \end{example}
}{% presentation
}{% nopresentation
}{% comments
}

\content{
        \begin{definition}
        In cluster sampling, the population is divided into subgroups
        (clusters), and a set of  subgroups are selected to be in the
        sample.
        \end{definition}
}{% presentation
\vspace{2in}
}{% nopresentation
\vspace{1.5in}
}{% comments
}

\content{
        \begin{example} If the college wanted to survey students,
        since students are already divided into classes, they could
        simply randomly select 50 classes and give the survey to all
        students in those classes. 
        \end{example}
}{% presentation
}{% nopresentation
}{% comments
}

\content{
        \begin{definition}
        In systematic sampling, every $n^{th}$ member of the population is
        selected to be in the sample.
        \end{definition}
}{% presentation
\vspace{2in}
}{% nopresentation
\vspace{1.5in}
}{% comments
}

\content{
        Perhaps the worst two sampling methods are the following two.
        \begin{definition}
        Convenience sampling is samples chosen by selecting whoever is
        convenient.
        \end{definition}
        \begin{definition}
        Voluntary response sampling is allowing the sample to volunteer.
        \end{definition}
        What are the potential problems with these two methods? 
}{% presentation
}{% nopresentation
}{% comments
}

\newpage
\subsection{How to mess things up before you start}


\content{
        We have already seen a number of sources of bias: 
        \begin{itemize}
        \item \textbf{Sampling bias} - when the sample is not
        representative of the population.
        \item \textbf{Voluntary response bias} - volunteers may not
        represent the target population
        \end{itemize}
}{% presentation
}{% nopresentation
}{% comments
}

\content{
        Sources of bias: 
        \begin{itemize}
        \item \textbf{Sampling bias} – when the sample is not representative of
        the population  
        \item \textbf{Voluntary response bias} – the sampling bias
        that often occurs when the sample is volunteers
        \item \textbf{Self-interest study} – bias that can occur when
        the researchers have an interest in the outcome 
        \item \textbf{Response bias} – when the
        responder gives inaccurate responses for any reason  
        \item \textbf{Perceived lack of anonymity} – 
        when the responder fears giving an honest
        answer  might negatively affect them  
        \item \textbf{Loaded questions} – when
        the question wording influences the responses  
        \item \textbf{Non-response bias} – when people refusing to
        participate in the study can influence the validity of the
        outcome.
        \end{itemize}
}{% presentation
}{% nopresentation
}{% comments
}

\content{
        \begin{example}
        Consider a recent study which found that chewing gum may raise
        math grades in teenagers. This study was conducted by the
        Wrigley Science Institute, a branch of the Wrigley chewing
        gum company. This is an example of a self-interest study; one
        in which the researches have  a vested interest in the outcome
        of the study. While this does not necessarily ensure that the
        study was biased, it certainly suggests that we should subject
        the study to extra scrutiny. 
        \end{example}
}{% presentation
\vspace{1in}
}{% nopresentation
\vspace{1in}
}{% comments
}

\content{
        \begin{example}
        A survey asks people “when was the last time you visited your
        doctor?”  
        \end{example}
}{% presentation
\vspace{1in}
}{% nopresentation
\vspace{1in}
}{% comments
}

\content{
        \begin{example}
        A survey asks “do you support funding research of alternative
        energy sources to reduce our  reliance on high-polluting
        fossil fuels?”  
        \end{example}
}{% presentation
\vspace{1in}
}{% nopresentation
\vspace{1in}
}{% comments
Loaded question. Make comment about order of questions as well.
}

\content{
        \begin{example} A researcher wishes to study drug-use in local
        high schools. They ask students how many times they used
        marijuana in the last 6 months.
        \end{example}
}{% presentation
\vspace{1in}
}{% nopresentation
\vspace{1in}
}{% comments
}
